% This file is part of the modular source for Stratified Relational Workspace Theory.
% !TeX root = ../Stratified_Relational_Workspace_Theory.tex
\theorypart{VII}{Relational Irreducibility}

\topic{22}{Coalition Information and Bayes Risk}

For \(A\subseteq V\), the coalition statistic is \(\mathbf H_t^A\). Given an
action space \(\mathsf A\) and loss \(\ell\), define its unrestricted Bayes
risk by
\[
\boxed{
\Risk_\ell(A)
=
\inf_d
\mathbb E\!\left[
\ell\bigl(Y_t^+,d(\mathbf H_t^A)\bigr)
\right].
}
\]
The infimum is over measurable decision rules for which the expectation is
defined. The full-system risk is \(\Risk_\ell(V)\).

\begin{proposition}[Coalition monotonicity]\label{prop:coalition-monotonicity}
If \(A\subseteq B\subseteq V\), then
\[
\Risk_\ell(B)\leq\Risk_\ell(A).
\]
\end{proposition}

\begin{proof}
Every decision rule available to coalition \(A\) is also available to coalition
\(B\), which can ignore its additional information.
\end{proof}

For finite \(\mathsf Y\), logarithmic loss converts risk gaps into conditional
information exactly.

\begin{theorem}[Coalition log-risk identity]\label{thm:coalition-log-risk}
Assume \(H(Y_t^+)<\infty\). If \(A\subseteq B\subseteq V\), then
\[
\boxed{
\Risk_{\log}(A)-\Risk_{\log}(B)
=
I\!\left(
Y_t^+;\mathbf H_t^B
\mid
\mathbf H_t^A
\right).
}
\]
In particular,
\[
\Risk_{\log}(A)-\Risk_{\log}(V)
=
I\!\left(
Y_t^+;\mathbf H_t
\mid
\mathbf H_t^A
\right).
\]
\end{theorem}

\begin{proof}
Under logarithmic loss the optimal report is the conditional target law, so
\[
\Risk_{\log}(A)=H(Y_t^+\mid\mathbf H_t^A).
\]
Because \(\mathbf H_t^A\) is a function of \(\mathbf H_t^B\), subtracting the
two conditional entropies gives the stated conditional mutual information.
\end{proof}

The theorem makes the decomposition explicit: irreducibility is not an
unqualified property of a system. It is a property relative to the chosen node
partition, accessible histories, target, and decision loss.

\topic{23}{Joint Advantage and Its Scope}

Define the coalition gap
\[
\Delta_A^{\ell}
=
\Risk_\ell(A)-\Risk_\ell(V)
\geq0.
\]
The node-level joint advantage is
\[
\boxed{
\Delta_{\mathrm{joint}}^{\ell}
=
\min_{i\in V}
\Delta_{\{i\}}^{\ell}.
}
\]
Under logarithmic loss,
\[
\boxed{
\Delta_{\mathrm{joint}}^{\log}
=
\min_{i\in V}
I\!\left(
Y_t^+;\mathbf H_t
\mid
\mathbf H_t^{\{i\}}
\right).
}
\]
Therefore
\[
\Delta_{\mathrm{joint}}^{\log}>0
\]
means that every individual component omits some predictive information
available to the complete system.

This quantity should not automatically be identified with a particular notion
of partial-information synergy. A positive gap can arise from complementary
information distributed across components, from higher-order interaction, or
from both. The theory uses the operationally weaker and more directly
testable statement:
\[
\text{no individual node is sufficient for the declared target.}
\]

For \(0\leq r\leq |V|\), define the order-\(r\) coalition gap profile by
\[
\boxed{
\Delta^{(r)}_\ell
=
\min_{\substack{A\subseteq V\\|A|\leq r}}
\left[
\Risk_\ell(A)-\Risk_\ell(V)
\right].
}
\]
The profile is nonincreasing in \(r\) and reaches zero by \(r=|V|\).
It records how rapidly operative capacity becomes reconstructible as larger
coalitions are permitted.

\topic{24}{Minimal Sufficient Coalition}

For tolerance \(\varepsilon\geq0\), define
\[
\boxed{
k_\varepsilon
=
\min\left\{
|A|:
\Risk_\ell(A)
\leq
\Risk_\ell(V)+\varepsilon
\right\}.
}
\]
The set of minimal coalitions is
\[
\mathcal K_\varepsilon
=
\left\{
A\subseteq V:
|A|=k_\varepsilon,
\Risk_\ell(A)
\leq
\Risk_\ell(V)+\varepsilon
\right\}.
\]

\begin{corollary}[Information form of coalition size]
Under the assumptions of Theorem~\ref{thm:coalition-log-risk},
\[
\boxed{
k_\varepsilon
=
\min\left\{
|A|:
I(Y_t^+;\mathbf H_t\mid\mathbf H_t^A)
\leq\varepsilon
\right\}.
}
\]
Equivalently, \(k_\varepsilon\) is the first \(r\) for which
\(\Delta_{\log}^{(r)}\leq\varepsilon\).
\end{corollary}

This yields a graded classification:
\[
\begin{aligned}
k_0=0
&\quad\Longrightarrow\quad
\text{the background law alone is sufficient},\\
k_0=1
&\quad\Longrightarrow\quad
\text{some individual node is sufficient},\\
k_0\geq2
&\quad\Longrightarrow\quad
\text{node-level relational irreducibility},\\
k_0=|V|
&\quad\Longrightarrow\quad
\text{every proper coalition is insufficient}.
\end{aligned}
\]
The last condition is much stronger and should not be inferred from the first.

\topic{25}{Reconstructibility of the Canonical State}

Let \(d_{\mathcal P}\) be a bounded metric generating the measurable structure
on \(\mathcal P(\mathsf Y)\). Define the reconstruction defect of the canonical
state from coalition \(A\) by
\[
\delta_A^Z
=
\inf_{\psi_A}
\mathbb E\!\left[
 d_{\mathcal P}\!\left(
 Z_t,\psi_A(\mathbf H_t^A)
 \right)
\right].
\]

\begin{proposition}[Sufficiency--reconstruction equivalence]
\label{prop:sufficiency-reconstruction}
For a coalition \(A\subseteq V\), the following are equivalent:
\begin{enumerate}
\item \(\mathbf H_t^A\) is sufficient for \(Y_t^+\) relative to the full
 history;
\item \(Z_t\) is measurable with respect to \(\mathcal F_t^A\);
\item there exists a measurable \(\psi_A\) such that
\[
Z_t=\psi_A(\mathbf H_t^A)
\qquad\text{almost surely}.
\]
\end{enumerate}
Each condition implies \(\delta_A^Z=0\). Under finite-target logarithmic loss,
they are also equivalent to
\[
\Risk_{\log}(A)=\Risk_{\log}(V).
\]
If the reconstruction infimum is known to be attained, then
\(\delta_A^Z=0\) is itself equivalent to the three conditions.
\end{proposition}

\begin{proof}
The equivalence of the first two statements is
Theorem~\ref{thm:predictive-state}; the second and third are two formulations
of measurability up to null sets. An exact reconstructor has zero expected
metric error. The log-risk statement follows from
Theorem~\ref{thm:coalition-log-risk} and the zero condition for conditional
mutual information. Without an attainment or compactness condition, an
infimum of zero alone certifies arbitrarily accurate reconstruction but need
not supply a single exact minimizing map.
\end{proof}

For an arbitrary enriched workspace variable \(W_t\), local
non-reconstructibility need not imply predictive irreducibility: \(W_t\) may
contain task-irrelevant material. The equivalence is exact for the canonical
state \(Z_t\), or for a candidate \(W_t\) known to be minimally aligned with
it.

\topic{26}{Matched Relational Nulls}

A null control must preserve the local facts that the claim is not intended to
challenge while destroying the relational organization being tested. Merely
reducing model capacity or replacing the workspace by noise is not a matched
null.

Let \(P\) be the intact law of
\((Y_t^+,\mathbf H_t^1,\ldots,\mathbf H_t^n)\). Define the local-marginal null
class
\[
\mathcal N_{\mathrm{marg}}(P)
=
\left\{
Q:
Q_{\mathbf H_t^i}=P_{\mathbf H_t^i}
\text{ for every }i
\right\},
\]
and the stronger local-predictive-matched class
\[
\boxed{
\mathcal N_{\mathrm{pred}}(P)
=
\left\{
Q:
Q_{(Y_t^+,\mathbf H_t^i)}
=
P_{(Y_t^+,\mathbf H_t^i)}
\text{ for every }i
\right\}.
}
\]
Because \(\mathbf H_t^i\) can include the complete local trajectory, the latter
class can preserve all single-component temporal and predictive behavior while
altering cross-component organization.

\begin{proposition}[Local-risk invariance under predictive matching]
\label{prop:local-risk-null}
If \(Q\in\mathcal N_{\mathrm{pred}}(P)\), then for every node \(i\) and every
loss determined by \((Y_t^+,d(\mathbf H_t^i))\),
\[
\boxed{
\Risk_{\ell,Q}(\{i\})
=
\Risk_{\ell,P}(\{i\}).
}
\]
\end{proposition}

\begin{proof}
Every such risk depends only on the joint law of
\((Y_t^+,\mathbf H_t^i)\), which is identical under \(P\) and \(Q\).
\end{proof}

A valid null must also satisfy a declared relation-breaking condition. Examples
include:
\begin{itemize}
\item permuting source--target message pairings while preserving each message
 marginal;
\item shuffling temporal correspondence across nodes while preserving every
 local trajectory;
\item permuting role bindings while preserving the local token or state
 distributions;
\item replacing an artifact by a matched-size, matched-format artifact whose
 relation to the surrounding history has been destroyed; and
\item randomizing transfer channels through a physical or software
 intervention while preserving message rate and marginal content.
\end{itemize}

For a predeclared null \(Q\), define
\[
\boxed{
\Delta_{\mathrm{tr}}(Q)
=
\Risk_{\ell,Q}(V)
-
\Risk_{\ell,P}(V).
}
\]
For a finite validated panel
\(\mathcal Q=\{Q_1,\ldots,Q_m\}\subseteq\mathcal N_{\mathrm{pred}}(P)\), report
both the complete vector
\[
\boldsymbol\Delta_{\mathrm{tr}}(\mathcal Q)
=
\bigl(
\Delta_{\mathrm{tr}}(Q_1),\ldots,
\Delta_{\mathrm{tr}}(Q_m)
\bigr)
\]
and the robust panel minimum
\[
\Delta_{\mathrm{tr}}^{\min}(\mathcal Q)
=
\min_{Q\in\mathcal Q}\Delta_{\mathrm{tr}}(Q).
\]
The panel must be fixed independently of the observed favorable gap. Otherwise
null selection can manufacture apparent relational dependence.

\topic{27}{Relational Witness Pairs}

\begin{definition}[Relational witness pair]\label{def:relational-witness}
A pair of laws \((P,Q)\) is a predictive relational witness pair for the
decomposition \(V\) and target \(Y_t^+\) when
\[
Q\in\mathcal N_{\mathrm{pred}}(P)
\]
but
\[
\Risk_{\ell,Q}(V)
\neq
\Risk_{\ell,P}(V).
\]
The two systems agree on every single-node predictive marginal but differ in
full-system operative capacity.
\end{definition}

\begin{example}[XOR witness pair]\label{ex:xor-witness}
Under the intact law \(P\), let
\[
X_1,X_2\overset{\mathrm{iid}}{\sim}\operatorname{Bernoulli}(1/2),
\qquad
Y=X_1\oplus X_2.
\]
Then each \((X_i,Y)\) pair is independent and uniform, so
\[
H_P(Y\mid X_i)=1\text{ bit},
\qquad
H_P(Y\mid X_1,X_2)=0.
\]
Let \(Q\) make \(X_1,X_2,Y\) mutually independent fair bits. Then
\[
Q_{(X_i,Y)}=P_{(X_i,Y)}
\qquad(i=1,2),
\]
but
\[
H_Q(Y\mid X_1,X_2)=1\text{ bit}.
\]
Therefore
\[
\Delta_{\mathrm{joint}}^{\log}=1\text{ bit},
\qquad
k_0=2,
\qquad
\Delta_{\mathrm{tr}}(Q)=1\text{ bit}.
\]
This proves predictive relational structure relative to the declared variables.
It does not, by itself, prove that any workspace variable is causally active.
\end{example}

\begin{example}[Redundant-copy control]
Let \(Y\) be a fair bit and set \(X_1=X_2=Y\). The two nodes are perfectly
coherent and the full system predicts perfectly, but
\[
H(Y\mid X_i)=0,
\qquad
\Delta_{\mathrm{joint}}^{\log}=0,
\qquad
k_0=1.
\]
Duplication is therefore not relational irreducibility.
\end{example}

The witness-pair criterion is stronger than a performance comparison against a
smaller model because it holds local predictive behavior fixed. It remains
observational. Part~VIII adds interventions and converts predictive
relationality into an operational causal claim.
